\section{Discussion}
\label{sec:discussion}

The central finding of this paper is that uncertainty estimation alone is
not enough. Without a deferral mechanism to act on the uncertainty signal,
the information it provides goes unused. The gap between ``having uncertainty''
and ``using it to reduce errors'' is the gap this work measures. TTA with
adaptive deferral reduces segmentation error by roughly 80\%; the same TTA
uncertainty without any deferral policy reduces error by 0\%---the predictions
are identical, only the decision changes. Choosing TTA over MC Dropout
improves Unc-AUROC by roughly 16 points in our experiments.
Choosing the policy well changes how much value is extracted from that
uncertainty. Neither
choice alone gets you to the best operating point; both matter, and they
interact. Figure~\ref{fig:summary_scatter} captures this visually: the six
method--policy configurations span a wide region of the deferral design space,
and the best position (upper-left: high error reduction, low deferral rate)
requires getting both choices right.

\subsection{Why TTA produces better uncertainty than MC Dropout}

The large Unc-AUROC gap between TTA (0.881) and MC Dropout (0.722)
is the largest effect in our experiments.

MC Dropout approximates a posterior over weights by sampling binary dropout
masks. The resulting uncertainty reflects how much the model's predictions
change when random subsets of neurons are disabled. This is a global
perturbation: it affects the entire network, and its effect on any given
pixel depends on how that pixel's representation propagates through all
layers. The result, empirically, is uncertainty that is spatially smooth.
Entire regions of the image receive similar uncertainty values because
nearby pixels share most of their computational path through the encoder.

TTA, by contrast, perturbs the input. A horizontal flip changes which
pixels fall on vessel boundaries; a 90-degree rotation changes the
orientation of thin vessels relative to the convolutional filters. The
resulting disagreement is spatially specific: it concentrates on pixels
whose predictions depend on local geometric context, which is exactly
where segmentation errors occur.

A subtler factor: MC Dropout uncertainty is computed
as mutual information (Equation~\ref{eq:mi}), which requires the
per-pass predictions $\hat{p}^{(t)}$ to disagree. But dropout perturbations
are often too small to flip a prediction: if the model is confident that
a pixel is background ($\hat{p} \approx 0.05$), 30 dropout passes might
produce values in $[0.03, 0.08]$, all on the same side of the decision
boundary. The mutual information is low even though the pixel might be
wrong. TTA augmentations, being geometric transformations, can produce
larger disagreements because they fundamentally change which input features
are aligned with which filter orientations.

Effective sample count also matters. With $T = 30$ MC Dropout
passes, the effective number of independent samples is lower than 30
because dropout masks overlap substantially (any two masks share roughly
$1 - 2 \times 0.3 + 0.3^2 = 0.49$ of their active neurons, assuming
$p = 0.3$). With $K = 6$ TTA augmentations, each augmentation is
genuinely different (identity, flip, rotate). The information content
per sample is higher.

\subsection{Why adaptive deferral outperforms global thresholding}

Global thresholding applies a single $\tau$ across all images. In a dataset
with variable image difficulty, this creates two failure modes: on easy
images, few pixels exceed $\tau$ and deferral contributes little (those
pixels were probably correct anyway); on hard images, many pixels exceed
$\tau$ and the system defers too aggressively, wasting review capacity on
pixels that are uncertain but correct.

Adaptive thresholding normalizes by image, deferring the top $\alpha$\%
most uncertain pixels regardless of absolute uncertainty magnitude. This
guarantees that each image contributes proportionally to the deferral
budget. The per-image percentile approach is also label-free at test time:
it requires no ground truth, only the uncertainty distribution of the
current image.

The confidence-aware strategy addresses a different failure mode. Both
global and adaptive thresholding defer based on uncertainty alone, but
uncertainty and error are not the same thing. A pixel can be uncertain
(moderate variance across passes) yet correctly predicted (mean probability
far from 0.5). The confidence-aware score $s = u \cdot (1 - c)$ suppresses
deferral for such pixels, concentrating the deferral budget on pixels
that are both uncertain and borderline. In our results, this produces
the best low-budget efficiency for TTA (4.56 for confidence-aware deferral,
versus 4.19 for the global TTA operating point at a slightly larger review
budget).

The adaptive and confidence-aware strategies are complementary. Adaptive
deferral works best at high deferral rates (25--30\%), where per-image
normalization prevents easy images from being ignored. Confidence-aware
deferral works best at low deferral rates (5--12\%), where concentrating on
the highest-risk pixels matters most. A hybrid policy, adaptive
thresholding with confidence weighting, is a natural extension we leave
for future work.

\subsection{Why calibration fails to improve deferral}

The mechanism is explained in Section~\ref{sec:calibration}: temperature
scaling is a monotonic transform that preserves uncertainty rankings. But
there is a deeper issue. ECE and deferral quality optimize genuinely
different objectives. ECE asks that errors be distributed proportionally
across confidence bins. Deferral asks that errors concentrate at high
uncertainty. A model that satisfies both simultaneously would need errors
to be both uniformly spread (for ECE) and concentrated (for deferral),
which is a contradiction except in trivial cases.

Laves et al.\ \citep{laves2020well} observed this tension for classification.
Our results confirm it holds for dense pixel-level prediction. The practical
takeaway is to report calibration metrics and decision metrics side by side:
Unc-AUROC, risk-coverage behavior, and achieved operating points answer a
different question from ECE.

\subsection{Deployment recommendations}

Based on the results in Sections~\ref{sec:results}--\ref{sec:ablations},
a practical deployment should:

\begin{enumerate}[leftmargin=*]
\item Use TTA (6 geometric transforms) for uncertainty when the goal is
pixel-level review triage. It dominates MC Dropout on both quality and speed
in our experiments and requires no architectural changes.

\item Use confidence-aware deferral for review budgets $\leq$ 15\% of pixels,
adaptive thresholding for larger budgets. Both are computed from the existing
prediction and uncertainty maps at negligible cost.

\item Skip temperature scaling for deferral. If calibrated probabilities are
needed for other purposes (e.g., reporting confidence), compute them
separately; the deferral policy should use uncalibrated uncertainty.

\item Monitor Unc-AUROC on a small labeled subset when deploying across sites.
If it drops below $\sim$0.70, switch from pixel-level deferral to image-level
flagging.

\item Match the deferral rate to the workflow: around 10--12\% if review
budget is tight and around 25\% if maximizing error removal matters more than
coverage.
\end{enumerate}
